%% =============================================================================
%%  pt.tex — textos da interface em português
%%  Carregado por estilo/estudobiblico.sty quando \ebIdioma é "pt".
%%  Para criar um idioma novo: copie este arquivo, traduza os textos e salve
%%  como <codigo>.tex (ex.: fr.tex). Depois use \newcommand\ebIdioma{fr}.
%%  NÃO mude as chaves (o primeiro argumento de \EBdeftxt).
%% =============================================================================

% Língua principal do babel (hifenização, datas, nomes automáticos)
\def\EB@babelmain{brazilian}
% Código da coluna de nomes dos livros bíblicos em dados/biblia/livros.tsv
% ("" = coluna "nome"; "en" = coluna "nome_en"; e assim por diante)
\def\EB@colunalivros{}

%% ---------------------------------------------------------------- estrutura
\EBdeftxt{numeroAbrev}{nº}
\EBdeftxt{etapa}{Etapa}
\EBdeftxt{parte}{Parte}
\EBdeftxt{sumario}{Sumário}
\EBdeftxt{bibliografia}{Bibliografia}
\EBdeftxt{indiceRemissivo}{Índice remissivo}
\EBdeftxt{indiceTermos}{Índice de termos originais}
\EBdeftxt{indiceTextos}{Índice de textos bíblicos}

%% ------------------------------------------------------------------- caixas
\EBdeftxt{colofaoUso}{Material de estudo pessoal.}
\EBdeftxt{colofaoFontes}{Texto grego: SBL Greek New Testament (CC BY 4.0, Society of Biblical Literature \& Logos); morfologia MorphGNT (CC BY-SA 3.0). Hebraico e aramaico: Westminster Leningrad Codex com morfologia da OpenScriptures Hebrew Bible (CC BY 4.0). Os textos das traduções modernas pertencem aos respectivos detentores de direitos e são citados para estudo.}
\EBdeftxt{colofaoTipografia}{Tipografia: EB Garamond, Gentium Book Plus, Marcellus, Source Sans 3, Noto Serif Hebrew e Noto Sans Syriac (SIL Open Font License). Composto em \LaTeX\ (LuaLaTeX).}

\EBdeftxt{objetivo}{Objetivo}
\EBdeftxt{nota}{Nota}
\EBdeftxt{pergunta}{Pergunta}
\EBdeftxt{cuidado}{Cuidado}
\EBdeftxt{aplicacao}{Aplicação}
\EBdeftxt{criticaTextual}{Crítica textual}
\EBdeftxt{baseAT}{Base AT:}
\EBdeftxt{baseNT}{Base NT:}
\EBdeftxt{traducoes}{Traduções}
\EBdeftxt{noNT}{no NT}

%% -------------------------------------------------------------- interlinear
\EBdeftxt{grego}{Grego}
\EBdeftxt{hebraico}{Hebraico}
\EBdeftxt{aramaico}{Aramaico}
\EBdeftxt{siriaco}{Siríaco}
\EBdeftxt{corpusNT}{No NT}
\EBdeftxt{corpusAT}{No AT}
\EBdeftxt{camadas}{Camadas}
\EBdeftxt{original}{original}
\EBdeftxt{transliteracaoMin}{transliteração}
\EBdeftxt{glosaPropria}{glosa própria}
\EBdeftxt{glosaLexical}{glosa lexical (en)}
\EBdeftxt{lemaMin}{lema}
\EBdeftxt{morfologiaMin}{morfologia}
\EBdeftxt{varianteMin}{variante}
\EBdeftxt{implicitoMin}{implícito}
\EBdeftxt{variantes}{Variantes}
\EBdeftxt{cabRef}{Ref.}
\EBdeftxt{cabPalavra}{Palavra}
\EBdeftxt{cabTransliteracao}{Transliteração}
\EBdeftxt{cabLema}{Lema}
\EBdeftxt{cabMorfologia}{Morfologia}
\EBdeftxt{cabGlosa}{Glosa}
\EBdeftxt{cabStrong}{Strong}
\EBdeftxt{cabNoTrecho}{No trecho}
\EBdeftxt{cabGlosasNoTrecho}{Glosas no trecho}

%% ---------------------------------------------------------- caderno de livro
\EBdeftxt{licao}{Lição}
\EBdeftxt{licoes}{Lições}
\EBdeftxt{seloLeia}{LEIA}
\EBdeftxt{citado}{citado}
\EBdeftxt{textosBiblicos}{Textos bíblicos}
\EBdeftxt{dicaTextos}{\S\ = ponto da lição · marque ao ler o contexto (capítulo) e ao comparar traduções}
\EBdeftxt{semTextos}{(nenhum texto bíblico encontrado nesta lição)}
\EBdeftxt{nTextos}{textos bíblicos}
\EBdeftxt{paraLer}{para ler}
\EBdeftxt{estudadaEm}{Estudada em}
\EBdeftxt{revisadaEm}{Revisada em}
\EBdeftxt{cabNum}{N\textsuperscript{o}}
\EBdeftxt{cabPag}{Pág.}
\EBdeftxt{cabTextos}{Textos}
\EBdeftxt{cabEstudada}{Estudada}
\EBdeftxt{cabRevisada}{Revisada}
\EBdeftxt{cabTipo}{Tipo}
\EBdeftxt{cabContexto}{Contexto}
\EBdeftxt{cabTraducoes}{Traduções}
\EBdeftxt{cabObservacao}{Observação}
\EBdeftxt{cabTexto}{Texto}
\EBdeftxt{cabVezes}{Vezes}
\EBdeftxt{cabEmQue}{Em que lições}
\EBdeftxt{antesDeLer}{Antes de ler}
\EBdeftxt{antesDeLerDica}{Que pergunta esta lição responde? (em minhas palavras)}
\EBdeftxt{ideiaCentral}{Ideia central}
\EBdeftxt{ideiaCentralDica}{a resposta da lição, em uma ou duas frases}
\EBdeftxt{anotacoes}{Anotações}
\EBdeftxt{anotacoesDica}{pontos principais de cada seção; o que me chamou a atenção}
\EBdeftxt{duvidas}{Dúvidas para pesquisar}
\EBdeftxt{revisaoSemConsultar}{Revisão sem consultar}
\EBdeftxt{revisaoSemConsultarDica}{responda às perguntas de revisão da lição de memória}
\EBdeftxt{escriturasHebraicas}{Escrituras Hebraicas}
\EBdeftxt{escriturasGregas}{Escrituras Gregas}

%% -------------------------------------------------------- caderno temático
\EBdeftxt{pesoCentral}{CENTRAL}
\EBdeftxt{pesoApoio}{apoio}
\EBdeftxt{pesoDificil}{DIFÍCIL}
\EBdeftxt{cabGrupo}{Grupo}
\EBdeftxt{cabPeso}{Peso}
\EBdeftxt{cabLido}{Lido}
\EBdeftxt{cabNota}{Nota}
\EBdeftxt{cabTranslitAbr}{Translit.}
\EBdeftxt{cabOcorr}{Ocorr.}
\EBdeftxt{cabSentido}{Sentido}
\EBdeftxt{siglaAT}{AT}
\EBdeftxt{siglaNT}{NT}
\EBdeftxt{parteDaBiblia}{Parte da Bíblia}
\EBdeftxt{textosEstudados}{Textos estudados}
\EBdeftxt{ocorrenciasPalavras}{Ocorrências das palavras}
\EBdeftxt{base}{Base:}
\EBdeftxt{temaVazio}{(nenhum texto com esse critério em textos.tsv)}
